
\documentclass[12pt,a4paper]{article}
\usepackage{amsmath}
\usepackage{amsfonts}
\usepackage{amssymb}
\usepackage{geometry}
\geometry{margin=1in}
\usepackage{hyperref}
\hypersetup{colorlinks=true, linkcolor=blue, urlcolor=cyan}

\title{Model Predictive Control (MPC) - Öğrenim ve Uygulama Raporu}
\date{\today}
\author{}

\begin{document}
\maketitle
\tableofcontents
\newpage

\section{Giriş}
Model Predictive Control (MPC), sistemin gelecekteki davranışını tahmin ederek en uygun kontrol sinyalini belirleyen ileri seviye bir kontrol yöntemidir. 
MPC, sistem kısıtlarını (örn: maksimum voltaj, maksimum hız) doğrudan hesaba katabilir ve özellikle elektrikli motorlar, otonom sistemler ve enerji yönetimi alanlarında yaygın olarak kullanılır.

\section{DC Motor Sistem Modeli}

\subsection{Durum Uzayı Temsili}
Elektriksel Dinamik:
\begin{equation}
V(t) = L \frac{di(t)}{dt} + Ri(t) + K_e \omega(t)
\end{equation}
Mekanik Dinamik:
\begin{equation}
J \frac{d\omega(t)}{dt} + B\omega(t) = K_t i(t)
\end{equation}

\subsection{Durum Değişkenleri}
\begin{itemize}
\item $x_1 = i(t)$ : Akım
\item $x_2 = \omega(t)$ : Açısal Hız
\end{itemize}

\subsection{Durum Uzayı Modeli}
\begin{equation}
\dot{x} = A x + B u
\end{equation}
\begin{equation}
y = C x + D u
\end{equation}
Matrisler:
\begin{equation}
A = \begin{bmatrix} -\frac{R}{L} & -\frac{K_e}{L} \\ \frac{K_t}{J} & -\frac{B}{J} \end{bmatrix}, \quad
B = \begin{bmatrix} \frac{1}{L} \\ 0 \end{bmatrix}, \quad
C = \begin{bmatrix} 0 & 1 \end{bmatrix}, \quad
D = 0
\end{equation}

Sistem MATLAB ortamında diskretize edilerek MPC tasarımı için hazırlandı.

\section{Temel 1-Adım Horizon MPC}

\subsection{Prediction Horizon}
1 adım.

\subsection{Cost Fonksiyonu}
\begin{equation}
J = (y_{k+1} - r)^2 + \lambda u_k^2
\end{equation}

\subsection{Özellikler}
\begin{itemize}
\item Her adımda sadece bir adım ötesi tahmin edildi.
\item Optimal kontrol sinyali analitik olarak çözüldü.
\item Kontrol ceza katsayısı (lambda) ayarlanarak sistem davranışı optimize edildi.
\end{itemize}

\subsection{Sonuçlar}
\begin{itemize}
\item Sistem stabil bir şekilde hedef hıza ulaştı.
\item Kontrol kuvvetinin etkisi doğrudan gözlemlendi.
\end{itemize}

\section{Multi-Step (5 Adım Horizon) MPC}

\subsection{Prediction Horizon}
5 adım.

\subsection{Optimization Yöntemi}
Quadratic Programming (QP) kullanıldı.

\subsection{Cost Fonksiyonu}
\begin{equation}
J = \sum_{i=1}^{5} \left( (y_{k+i} - r)^2 + \lambda (u_{k+i-1})^2 \right)
\end{equation}

\subsection{Özellikler}
\begin{itemize}
\item Gelecek 5 adım boyunca sistemin davranışı tahmin edildi.
\item Optimal voltaj dizisi hesaplandı.
\item Sadece ilk kontrol sinyali uygulandı (receding horizon mantığı).
\end{itemize}

\subsection{Sonuçlar}
\begin{itemize}
\item Sistem hedefe daha hızlı ve optimize edilmiş bir yol izleyerek ulaştı.
\item Voltaj kısıtları ([-5V, +5V]) başarıyla uygulandı.
\item Overshoot ve dalgalanma gözlenmedi.
\end{itemize}

\section{Sonuçlar ve Gözlemler}
\begin{itemize}
\item 1 adım MPC basit ve doğrudan, ancak çok akıllı kararlar veremeyen bir yapı sundu.
\item Multi-Step MPC, daha planlı ve geleceği optimize ederek sistem davranışını iyileştirdi.
\item Kontrol ceza katsayısı (lambda) sistemin agresiflik veya sakinliği üzerinde doğrudan etkili oldu.
\end{itemize}

\section{Kazanımlar ve Gelecek Adımlar}

\subsection{Kazanımlar}
\begin{itemize}
\item MPC'nin temel mekanizması ve tahmin modeli anlaşıldı.
\item Diskretize sistem kullanarak MPC kurulumu öğrenildi.
\item Quadratic Programming ile optimal kontrol hesaplandı.
\item Sistem kısıtlarını MPC içine entegre etme pratiği kazanıldı.
\end{itemize}

\subsection{Potansiyel Gelişimler}
\begin{itemize}
\item Prediction horizon artırılarak daha akıllı kontrol.
\item Soft constraints eklenerek daha esnek sistem yönetimi.
\item Multi-Input Multi-Output (MIMO) sistemlerde MPC uygulaması.
\item Realtime MPC simülasyonları ve donanım üzerinde uygulamalar.
\item Neural Network destekli öğrenen MPC geliştirilmesi.
\end{itemize}

\end{document}
